% \documentclass[compacto,10pt]{aleph-notas}
\documentclass[compacto,10pt,comentarios]{aleph-notas}

% Se recomienda leer la documentación de esta
% clase en https://github.com/alephsub0/LaTeX_aleph-notas

% -- Paquetes adicionales
\usepackage{enumitem}
\usepackage{booktabs}
\usepackage{aleph-comandos}
\emergencystretch=2em


% -- Datos del libro
\institucion{Universidad Autónoma Metropolintna}
% \carrera{}
\asignatura{Diseño de experimentos en Ingeniería}
\tema{Aprende Diseño de Experimentos desde la práctica y la vida cotidiana}
\autor{Roman A. Mora}
\fecha{\today}
\fuente{montserrat}

%% --> Logos de las guias
\definecolor{colordef}{cmyk}{0.81,0.62,0.00,0.22}

\begin{document}

\encabezado

\section{¿Miel, azúcar o café?}

Actualmente muchas personas buscan alternativas más orgánicas para mejorar el crecimiento de las plantas. En algunos huertos caseros y sistemas de cultivo artesanal es común escuchar ideas como:

\begin{quote}
“La miel puede ayudar porque tiene nutrientes.”
\end{quote}

o bien:

\begin{quote}
“El café podría favorecer el crecimiento.”
\end{quote}

Incluso algunas personas consideran que el azúcar puede aportar energía al sistema.

Pero surge una pregunta importante:

\begin{center}
\textit{¿Realmente alguna de estas sustancias podría favorecer el crecimiento de una planta?}
\end{center}

y todavía más interesante:

\begin{center}
\textit{¿Todas producen el mismo efecto?}
\end{center}

Cuando aparecen varias posibilidades al mismo tiempo, comparar únicamente observando “a simple vista” puede resultar engañoso, ya que algunas diferencias podrían deberse solamente al azar.

Por esta razón construiremos un experimento controlado utilizando semillas de lenteja y diferentes tratamientos.

\section{Idea experimental}

En este experimento se evaluará el efecto de:

\begin{itemize}
\item miel,
\item azúcar,
\item y café soluble,
\end{itemize}

sobre la germinación y crecimiento inicial de semillas de lenteja.

Para ello se utilizará un medio semisólido elaborado con grenetina y se compararán los resultados obtenidos entre distintos grupos.

La idea no es únicamente observar si una lenteja germina o no, sino intentar responder preguntas como:

\begin{itemize}
\item ¿Cuál tratamiento produce mayor crecimiento?
\item ¿Cuál parece favorecer más la germinación?
\item ¿Cuál genera más humedad?
\item ¿Cuál presenta más contaminación?
\item ¿Existen diferencias visibles entre tratamientos?
\end{itemize}

\section{¿Por qué usar ANOVA?}

Imagine que después de varios días observa que algunas lentejas crecieron más en un tratamiento que en otro.

Entonces surge una pregunta fundamental:

\begin{center}
\textit{¿Las diferencias observadas realmente se deben al tratamiento o podrían deberse al azar?}
\end{center}

Aquí aparece una de las herramientas más importantes del diseño experimental:

\begin{center}
\textbf{ANOVA}
\end{center}

ANOVA significa:

\begin{center}
\textit{Análisis de Varianza}
\end{center}

y permite comparar grupos para determinar si las diferencias observadas entre ellos parecen suficientemente grandes como para pensar que existe un efecto real asociado al tratamiento aplicado.

En otras palabras:

\begin{quote}
ANOVA nos ayuda a comparar varios grupos al mismo tiempo.
\end{quote}

\section{Objetivo}

Evaluar el efecto de distintas sustancias sobre la germinación y crecimiento inicial de semillas de lenteja utilizando un diseño experimental basado en ANOVA de un factor.

\section{Materiales}

\begin{itemize}
\item Agua
\item Grenetina
\item Azúcar
\item Café soluble
\item Miel
\item Agua oxigenada comercial
\item Antihongos comercial
\item Semillas de lenteja
\item Recipientes plásticos
\item Pipeta
\item Báscula
\item Fuente de calor
\end{itemize}

\section{Construcción del medio de cultivo}

Para construir el medio semisólido se utilizaron:

\begin{table}[H]
\centering
\begin{tabular}{lc}
\toprule
Componente & Cantidad \\
\midrule
Agua & 2000 g \\
Grenetina & 100 g \\
Azúcar & 100 g \\
Agua oxigenada comercial & 30 g \\
Antihongos comercial & 0.5 g \\
\bottomrule
\end{tabular}
\caption{Composición del medio de cultivo}
\end{table}

\section{Preparación del medio}

\begin{enumerate}
\item Colocar el agua en un recipiente resistente al calor.
\item Agregar la grenetina y el azúcar.
\item Llevar el sistema a ebullición.
\item Mantener el hervor durante cinco minutos.
\item Retirar del calor.
\item Esperar una ligera disminución de temperatura.
\item Agregar posteriormente el agua oxigenada y el antihongos.
\item Mezclar cuidadosamente.
\item Verter aproximadamente 30 g de medio en cada recipiente.
\item Permitir la solidificación.
\end{enumerate}

\section{Recipientes utilizados}

Los recipientes utilizados poseen dimensiones aproximadas de:

\[
4 \times 4 \times 5 \text{ cm}
\]

En cada recipiente se colocaron aproximadamente:

\[
30 \text{ g de medio}
\]

lo que generó una altura cercana a:

\[
1 \text{ cm}
\]

de medio semisólido.

\section{Preparación de tratamientos}

Cada solución fue preparada utilizando:

\[
30 \text{ g de agua}
\]

y

\[
30 \text{ g de soluto}
\]

Los tratamientos utilizados fueron:

\begin{itemize}
\item Azúcar
\item Café soluble
\item Miel
\end{itemize}

\section{Preparación de soluciones}

\begin{enumerate}
\item Colocar 30 g de agua en un recipiente.
\item Llevar el agua a ebullición.
\item Agregar el soluto correspondiente.
\item Mezclar continuamente.
\item Mantener el sistema en reposo durante 30 minutos.
\item Permitir la precipitación del exceso de soluto.
\item Tomar cuidadosamente el sobrenadante sin agitar.
\end{enumerate}

\section{Tratamiento blanco}

Además de los tratamientos anteriores, se preparó un grupo blanco que únicamente contiene el medio de cultivo sin sustancias adicionales.

Este grupo permitirá comparar el crecimiento natural de las semillas.

\section{Distribución de tratamientos}

\begin{table}[H]
\centering
\begin{tabular}{lc}
\toprule
Color del envase & Tratamiento \\
\midrule
Amarillo & Blanco \\
Morado & Miel \\
Azul & Azúcar \\
Verde & Café soluble \\
\bottomrule
\end{tabular}
\caption{Distribución experimental}
\end{table}

\section{Preparación de semillas}

Las semillas de lenteja fueron hidratadas previamente con el objetivo de:

\begin{itemize}
\item favorecer la germinación,
\item disminuir variabilidad experimental,
\item e identificar semillas posiblemente inviables.
\end{itemize}

\section{Aplicación experimental}

\begin{enumerate}
\item Colocar seis semillas de lenteja sobre cada recipiente.
\item Aplicar tres gotas de la solución correspondiente utilizando una pipeta.
\item Mantener condiciones similares de iluminación y temperatura.
\item Realizar observaciones diarias durante siete días.
\end{enumerate}

\section{Réplicas experimentales}

Cada tratamiento será realizado por triplicado.

La unidad experimental corresponde al recipiente y no a cada semilla individual.

\section{Condiciones ambientales}

Los recipientes serán colocados en un sitio con:

\begin{itemize}
\item temperatura relativamente estable,
\item luz indirecta,
\item baja perturbación física,
\item y humedad moderada.
\end{itemize}

Los recipientes permanecerán cerrados utilizando sus tapas correspondientes.

Dos veces al día:

\begin{itemize}
\item por la mañana,
\item y por la noche,
\end{itemize}

serán abiertos aproximadamente durante:

\[
20 \text{ minutos}
\]

para permitir intercambio gaseoso y disminuir acumulación excesiva de humedad.

\section{Variables de respuesta}

Durante el experimento se registrarán diariamente las siguientes variables:

\begin{itemize}
\item Tiempo de germinación
\item Número de semillas germinadas
\item Altura del tallo
\item Longitud de raíz
\item Presencia de hongos
\item Cambios de coloración
\item Presencia de condensación
\end{itemize}

\section{Diseño experimental}

En este experimento:

\begin{itemize}
\item el factor corresponde al tipo de tratamiento,
\item los niveles son miel, azúcar, café y blanco,
\item y el diseño utilizado corresponde a un ANOVA de un factor.
\end{itemize}

\section{Tabla de registro experimental}

\begin{table}[H]
\centering
\scriptsize
\resizebox{\linewidth}{!}{%
\begin{tabular}{ccccccccc}
\toprule
Tratamiento & Repetición & Día & Germinación & Semillas germinadas & Altura del tallo & Longitud de raíz & Hongos & Coloración/Condensación \\
\midrule
&&&&&&&&\\
&&&&&&&&\\
&&&&&&&&\\
&&&&&&&&\\
&&&&&&&&\\
&&&&&&&&\\
&&&&&&&&\\
&&&&&&&&\\
&&&&&&&&\\
&&&&&&&&\\
&&&&&&&&\\
&&&&&&&&\\
&&&&&&&&\\
&&&&&&&&\\
&&&&&&&&\\
\bottomrule
\end{tabular}%
}
\caption{Formato sugerido para el registro experimental}
\end{table}




\end{document}